\documentclass[11pt]{article}

\usepackage[a4paper,margin=28mm]{geometry}
\usepackage[T1]{fontenc}
\usepackage{lmodern}
\usepackage{microtype}
\usepackage{amsmath,amssymb,amsthm,mathtools}
\usepackage{booktabs}
\usepackage{array}
\usepackage{enumitem}
\usepackage[hidelinks]{hyperref}
\setlength{\emergencystretch}{2em}

\newtheorem{theorem}{Theorem}[section]
\newtheorem{proposition}[theorem]{Proposition}
\newtheorem{corollary}[theorem]{Corollary}
\newtheorem{lemma}[theorem]{Lemma}
\theoremstyle{definition}
\newtheorem{definition}[theorem]{Definition}
\theoremstyle{remark}
\newtheorem{remark}[theorem]{Remark}

\newcommand{\PG}{\operatorname{PG}}
\newcommand{\F}{\mathbb{F}}
\newcommand{\cC}{\mathcal{C}}
\newcommand{\cE}{\mathcal{E}}
\newcommand{\cQ}{\mathcal{Q}}
\newcommand{\cF}{\mathcal{F}}
\newcommand{\Fix}{\operatorname{Fix}}
\newcommand{\Npair}{N_{\mathrm{pair}}}
\newcommand{\pospart}[1]{\left(#1\right)_{+}}

\title{Frobenius-equivariant pair extension of eight-arcs in
$\PG(2,25)$}
\author{Tavis Rudd}
\date{July 2026}

\begin{document}
\maketitle

\begin{abstract}
Let $\phi$ be the quadratic Frobenius involution of $\PG(2,s^2)$.  We give
an exact carrierwise count for fresh conjugate pairs
$\{P,\phi(P)\}$ that extend a $\phi$-invariant arc.  The count combines the
exact number of empty Baer lines with a charge from noninvariant secant
orbits.  An exact correction separates secant orbits whose fixed centers
make them invisible on a carrier from genuine collisions between visible
charges, and it classifies equality and every first-order excess level.

We then prove that every $\phi$-invariant eight-arc in $\PG(2,25)$ admits a
fresh conjugate-pair extension.  The zero- and four-fixed-point profiles
follow from incidence and moment arguments.  The exceptional
two-fixed-point profile is reduced, by projective normalization, to a
finite certificate checked by the Lean kernel.  The full theorem, the
semantic global pair count, and the profile split are formalized in Lean
without unproved declarations, custom axioms, or native-code evaluation.
Together with the general count, the order-five result covers every
prime-power base-field order $s\ge5$.
\end{abstract}

\section{Introduction}

An arc in a projective plane is a set of points no three of which are
collinear.  Ordinary completeness asks whether one further point can be
adjoined.  If the arc is preserved by a field automorphism, the natural
equivariant continuation unit is instead a complete Galois orbit.  Over a
quadratic extension this means either a fixed point or a nonfixed conjugate
pair.

Let $F=\F_s$, let $K=\F_{s^2}$, and let
\[
  \phi\colon [x:y:z]\longmapsto[x^s:y^s:z^s]
\]
act on $\PG(2,K)$.  The fixed points and fixed lines form the embedded Baer
subplane $\PG(2,F)$.  This paper studies the following question: when must a
$\phi$-invariant arc admit an extension by a fresh pair
$\{P,\phi(P)\}$?

The classical ingredients are elementary: point and line counts in a
projective plane, the fixed Baer subplane, Hilbert--90 normalization, and
secant counting; see Hirschfeld~\cite{Hirschfeld1998} and
Martin~\cite{Martin1967}.  The contribution here is their exact assembly
into an orbit-valued extension count.  The main results are:

\begin{enumerate}[label=(\roman*),leftmargin=2.1em]
\item an explicit lower bound for the number of legal conjugate-pair
      extensions of an arbitrary invariant arc;
\item an exact linewise correction that separates invisible secant centers
      from charge collisions, together with equality and excess criteria;
\item a uniform theorem that every invariant eight-arc in $\PG(2,25)$
      extends by a fresh conjugate pair and, together with the general
      estimate, covers every prime-power base order $s\ge5$; and
\item a square-root lower bound for arcs that are saturated with respect to
      conjugate-pair insertion.
\end{enumerate}

The last constant is not new: it is the classical Lunelli--Sce
$\sqrt{2q}$ scale~\cite{LunelliSce1964,BallLavrauw2019}, now obtained under
the weaker hypothesis of no equivariant pair extension.  We make no
sharpness claim.

\subsection*{Prior art and claim boundary}

Baker and Wantz~\cite{BakerWantz2005} use Frobenius-invariant arcs in the
Hughes plane and, in a maximality argument, consider a point together with
its Frobenius image.  This is genuine precedence for the conjugate-point
maneuver.  Their result does not count legal conjugate pairs on empty Baer
lines.  Jungnickel's arcs in Baer-semiplane complements
\cite{Jungnickel1987} and Ueberberg's geometries of Frobenius collineations
\cite{Ueberberg1997} are also adjacent but concern different incidence or
collineation problems.

Ordinary arc and MDS lengthening results
\cite{Alderson2005,AldersonBruenSilverman2007,Alderson2026} impose no
quadratic Galois-pair constraint.  Likewise, the classification of
Galois-invariant subsets of $\mathbf P^1$~\cite{LopezMaisnerNartXarles2002}
does not treat arbitrary plane arcs or their pair extensions.  Bounded
zbMATH Open, OpenAlex, Crossref, and source-level searches found no exact
precursor for the quantitative criterion below or for the uniform
$\PG(2,25)$ theorem.  This is negative search evidence, not a claim of
historical priority.

\section{Quadratic Frobenius and mate lines}

Let $\cC\subseteq\PG(2,K)$ be a $\phi$-invariant $k$-arc.  Write
\[
  f=|\cC\cap\PG(2,F)|,
  \qquad k=f+2e,
\]
so that $\cC$ is the disjoint union of $f$ fixed points and $e$ nonfixed
conjugate pairs.

For every nonfixed point $P$, the line $P\phi(P)$ is fixed by $\phi$; we
call it the \emph{mate line} of the pair $\{P,\phi(P)\}$.  Conversely, every
nonfixed conjugate pair has a unique mate line.  A fixed line contains
$s^2+1$ points, of which $s+1$ are fixed, and therefore contains exactly
\[
  N=\frac{s^2-s}{2}
\]
nonfixed conjugate pairs.

\begin{definition}\label{def:legal-pair}
A \emph{legal conjugate-pair extension} of $\cC$ is an unordered pair
$Q=\{P,\phi(P)\}$ such that $P\ne\phi(P)$, $Q\cap\cC=\varnothing$, and
$\cC\cup Q$ is an arc.  Let $\Npair(\cC)$ be the number of such pairs.
\end{definition}

\begin{lemma}[Pair-extension test]\label{lem:pair-test}
Let $P\ne\phi(P)$ and suppose neither point belongs to $\cC$.  Then
$\cC\cup\{P,\phi(P)\}$ fails to be an arc if and only if at least one of the
following holds:
\begin{enumerate}[label=(\alph*),leftmargin=2em]
\item $P$ lies on a secant of $\cC$;
\item $\phi(P)$ lies on a secant of $\cC$;
\item the mate line $P\phi(P)$ contains a point of $\cC$.
\end{enumerate}
\end{lemma}

\begin{proof}
Every collinear triple in the enlarged set either contains one new point
and two old points, giving (a) or (b), or contains both new points and one
old point, giving (c).  The converse is immediate.
\end{proof}

The useful carriers are therefore the fixed lines disjoint from $\cC$.
Let $\cE$ be this set of empty fixed lines.

\begin{lemma}[Empty fixed lines]\label{lem:empty-lines}
The number of empty fixed lines is
\begin{equation}\label{eq:E}
 E=|\cE|=s^2+s+1-\left(f(s+1)-\binom f2+e\right).
\end{equation}
\end{lemma}

\begin{proof}
The $f$ fixed selected points occupy
$f(s+1)-\binom f2$ fixed lines.  Inclusion--exclusion is exact because no
three selected fixed points are collinear.  Each selected conjugate pair
occupies its mate line.  These $e$ lines are distinct and contain no fixed
selected point, since $\cC$ is an arc.  Subtract from the $s^2+s+1$ fixed
lines.
\end{proof}

Among the $\binom k2$ old secants, exactly
\[
 I=\binom f2+e
\]
are fixed: those joining two fixed selected points and the $e$ mate lines.
The remaining secants occur in two-element Frobenius orbits.  Their number
is
\begin{equation}\label{eq:M}
 M=\frac{\binom k2-I}{2}=fe+e(e-1).
\end{equation}

\section{The quantitative pair-extension criterion}

\begin{theorem}[Quadratic pair-extension criterion]
\label{thm:pair-criterion}
Let $\cC$ be a $\phi$-invariant $k$-arc in $\PG(2,s^2)$ with profile
$k=f+2e$.  With $E$, $N$, and $M$ as in
\eqref{eq:E} and \eqref{eq:M},
\begin{equation}\label{eq:uniform-lower}
  \Npair(\cC)\ge E\,\pospart{N-M}
  =E\,\pospart{\frac{s^2-s}{2}-fe-e(e-1)}.
\end{equation}
In particular, if $E>0$ and $M<N$, then $\cC$ admits a legal conjugate-pair
extension.
\end{theorem}

\begin{proof}
Fix $\ell\in\cE$.  It carries exactly $N$ candidate conjugate pairs.  A
fixed old secant meets $\ell$ in a fixed point and destroys no nonfixed
candidate.  Consider a nonfixed old secant orbit $\{m,\phi(m)\}$.  Unless
$m\cap\ell$ is fixed, the two intersections $m\cap\ell$ and
$\phi(m)\cap\ell$ form one conjugate candidate pair; thus the orbit destroys
at most one candidate.  There are $M$ such secant orbits, so at least
$\pospart{N-M}$ candidates survive on $\ell$.

A survivor lies on no old secant, and its empty mate line contains no point
of $\cC$.  Lemma~\ref{lem:pair-test} makes it legal.  Candidate sets on
distinct fixed lines are disjoint because a nonfixed pair has a unique mate
line.  Summing over the $E$ carriers proves \eqref{eq:uniform-lower}.
\end{proof}

The preceding proof deliberately uses only the uniform upper bound $M$ on
the number of forbidden candidates per carrier.  Retaining the actual
forbidden support gives both a sharper bound and an exact identity.

For $\ell\in\cE$, let $\cQ_\ell$ be its $N$ candidate pairs and let
$\cF_\ell\subseteq\cQ_\ell$ be the pairs hit by at least one nonfixed old
secant orbit.  Then
\begin{equation}\label{eq:heterogeneous}
  \Npair(\cC)=\sum_{\ell\in\cE}
  \left(|\cQ_\ell|-|\cF_\ell|\right).
\end{equation}
The same formula applies when carrier capacities vary.  In particular, if
$|\cF_\ell|\le U_\ell$, then
\[
  \Npair(\cC)\ge
  \sum_{\ell\in\cE}\pospart{|\cQ_\ell|-U_\ell}.
\]

\section{Invisible centers and collision correction}

Fix $\ell\in\cE$.  Every nonfixed secant orbit has a fixed center
$m\cap\phi(m)$.  Let $B_\ell$ be the number of such orbits whose center
lies on $\ell$.  Those orbits meet $\ell$ only in a fixed point and are
invisible to its nonfixed candidates.

For $Q\in\cQ_\ell$, let $\mu_\ell(Q)$ be the number of visible secant
orbits that charge $Q$, and put
\[
 A_\ell=\sum_{Q\in\cQ_\ell}\mu_\ell(Q),
 \qquad
 R_\ell=\sum_{Q\in\cQ_\ell}
     \pospart{\mu_\ell(Q)-1}.
\]
Here $R_\ell$ is the redundancy caused by several visible obstruction
orbits hitting the same candidate.

\begin{proposition}[Exact linewise correction]\label{prop:linewise}
For every empty fixed line $\ell$,
\begin{align}
 A_\ell&=M-B_\ell,\label{eq:visible-mass}\\
 |\cF_\ell|&=A_\ell-R_\ell,\label{eq:support}\\
 |\cQ_\ell\setminus\cF_\ell|+M
   &=N+B_\ell+R_\ell.\label{eq:linewise-balance}
\end{align}
Consequently the first-order identity
$|\cQ_\ell\setminus\cF_\ell|+M=N$ holds exactly when every secant orbit is
visible on $\ell$ and the visible charge map is injective.
\end{proposition}

\begin{proof}
Exactly $B_\ell$ of the $M$ secant orbits are invisible, proving
\eqref{eq:visible-mass}.  The support of the multiplicity function
$\mu_\ell$ is $\cF_\ell$, and for every positive integer $a$ one has
$1=a-(a-1)$.  Summing this equality over the support gives
\eqref{eq:support}.  Substitution into
$|\cQ_\ell\setminus\cF_\ell|=N-|\cF_\ell|$ gives
\eqref{eq:linewise-balance}.  Equality in the first-order identity is
equivalent to $B_\ell=R_\ell=0$.
\end{proof}

Let
\[
 L=\Npair(\cC),\qquad
 B=\sum_{\ell\in\cE}B_\ell,
 \qquad
 R=\sum_{\ell\in\cE}R_\ell.
\]

\begin{corollary}[Aggregate equality and excess]
\label{cor:aggregate}
The exact aggregate balance is
\begin{equation}\label{eq:aggregate}
  L+EM=EN+B+R.
\end{equation}
Thus $L+EM=EN$ if and only if every secant-orbit center avoids every empty
fixed carrier and every local visible charge is injective.  More generally,
for each integer $t\ge0$,
\[
  L+EM=EN+t\quad\Longleftrightarrow\quad B+R=t.
\]
\end{corollary}

This is an algebraic equality and excess classification.  It does not
classify the invariant arcs for which $L$ is small, and no such structural
inverse theorem is claimed.

\section{Eight-arcs over \texorpdfstring{$\F_{25}$}{F25}}

We now take $s=5$, so each empty fixed line carries $N=10$ candidate
pairs.  The profile of an invariant eight-arc is
\[
 (f,e)\in\{(0,4),(2,3),(4,2),(6,1),(8,0)\}.
\]

Table~\ref{tab:q25-profiles} records the numerical data and the proved
lower bound for each profile.  In the exceptional row the bound $L\ge1$
records the kernel-checked existence theorem; the stronger observed census
minimum in Remark~\ref{rem:census} is not used.

\begin{table}[ht]
\centering
\small
\begin{tabular}{ccccc l}
\toprule
$f$ & $e$ & $E$ & $M$ & proved $L\ge$ & method \\
\midrule
$0$ & $4$ & $27$ & $12$ & $5$   & secant-index moments \\
$2$ & $3$ & $17$ & $12$ & $1$   & kernel-checked certificate \\
$4$ & $2$ & $11$ & $10$ & $4$   & fixed-center incidence \\
$6$ & $1$ & $9$  & $6$  & $36$  & uniform criterion \\
$8$ & $0$ & $11$ & $0$  & $110$ & uniform criterion \\
\bottomrule
\end{tabular}
\caption{The five profiles of invariant eight-arcs in $\PG(2,25)$.
Here $L=\Npair(\cC)$.}
\label{tab:q25-profiles}
\end{table}

The following simple center estimate is useful.  A secant orbit joining
two distinct selected conjugate pairs has a fixed center that lies on
neither participating mate line.  At most $f$ fixed-point star lines and
$e-2$ other mate lines through this center are occupied.  Since a fixed
point lies on $s+1$ fixed lines, each such cross-pair secant orbit is
invisible on at least
\begin{equation}\label{eq:cross-invisible}
  s+3-f-e
\end{equation}
empty carriers.

\begin{proposition}[The four nonexceptional profiles]
\label{prop:four-profiles}
Let $\cC$ be an invariant eight-arc in $\PG(2,25)$.
\begin{enumerate}[label=(\alph*),leftmargin=2em]
\item If $f=8$ or $f=6$, then $L\ge110$ or $L\ge36$, respectively.
\item If $f=4$, then $L\ge4$.
\item If $f=0$, then $L\ge5$.
\end{enumerate}
\end{proposition}

\begin{proof}
For $(f,e)=(8,0)$ one has $(E,M)=(11,0)$, and for $(6,1)$ one has
$(E,M)=(9,6)$.  Theorem~\ref{thm:pair-criterion} gives $110$ and $36$.

For $(f,e)=(4,2)$ one has $(E,M,N)=(11,10,10)$.  There are two cross-pair
secant orbits, each invisible on at least two empty carriers by
\eqref{eq:cross-invisible}.  Hence $B\ge4$, and
\eqref{eq:aggregate} gives $L=B+R\ge4$.

It remains to treat $(f,e)=(0,4)$.  Here
\[
  (E,M,N)=(27,12,10).
\]
All twelve nonfixed secant orbits are cross-pair orbits, and
\eqref{eq:cross-invisible} gives $B\ge48$.

We recall the second secant-index moment for an eight-arc:
\begin{equation}\label{eq:second-moment}
  \sum_{x\notin\cC}\binom{r_x}{2}=3\binom84=210,
\end{equation}
where $r_x$ is the number of secants of $\cC$ through $x$.  At fixed
external points, the first secant-index sum is $48$: the four fixed mate
secants contain six fixed points each, contributing $24$ incidences, and
the other $24$ secants each contain their unique fixed orbit center.  Since $r_x\le4$ and
$2\binom n2\le3n$ for $0\le n\le4$, their contribution to
\eqref{eq:second-moment} is at most $72$.

The four occupied fixed carriers are precisely the selected mate lines.
For an external nonfixed point $x$ on one of them, write $r_x=1+u_x$, where
the first secant is the mate line.  Every nonfixed secant meets at most the
two selected mate lines not containing its endpoints, so
$\sum u_x\le48$.  Since $u_x\le3$ and
$\binom{1+u}{2}\le2u$, the occupied-carrier contribution is at most $96$.
Thus the two endpoints of candidates on empty carriers contribute at least
$210-72-96=42$ to \eqref{eq:second-moment}.  Conjugate endpoints contribute
equally, and therefore
\[
  T:=\sum_{\ell\in\cE}\sum_{Q\in\cQ_\ell}
       \binom{\mu_\ell(Q)}2\ge21.
\]
Because $\mu_\ell(Q)\le4$ and
$\binom n2\le2(n-1)$ for $1\le n\le4$, we have $T\le2R$, whence $R\ge11$.
Finally, \eqref{eq:aggregate} gives
\[
  L=27(10-12)+B+R\ge-54+48+11=5.
\]
\end{proof}

The remaining profile $(f,e)=(2,3)$ has $(E,M,N)=(17,12,10)$.  The center
estimate gives $B\ge18$, but the aggregate identity would still require a
further collision estimate.  We close this profile by a finite reduction.

\begin{proposition}[Exceptional profile]\label{prop:f2}
Every invariant eight-arc in $\PG(2,25)$ with exactly two fixed selected
points has a fresh legal conjugate-pair extension.
\end{proposition}

\begin{proof}[Computer-assisted proof, kernel checked]
Represent $\F_{25}$ as $\F_5[\omega]/(\omega^2-2)$.  A base-field
projectivity sends the two fixed selected points to a prescribed ordered
pair.  Their stabilizer then sends the first admissible selected conjugate
pair to
\[
 [1:\omega:\omega],\qquad[1:-\omega:-\omega].
\]
Both transformations commute with Frobenius and preserve the arc and
pair-extension predicates.

After these reductions, the remaining finite slice has $46{,}056$ rows.
Of these, $39{,}012$ carry a checked witness that the proposed old set is
not an arc, while $7{,}044$ carry a checked fresh legal-pair witness.  The
rows are split into $1{,}639$ leaf modules and $303$ aggregate modules.
Each leaf is evaluated by Lean's kernel reduction.  The projective
transport theorem maps the checked survivor back to the original arc and
states explicitly that both new points are fresh and that the enlarged set
is an arc.  Section~\ref{sec:formalization} records the trust boundary.
\end{proof}

\begin{theorem}[Uniform order-five extension]\label{thm:q25}
Every Frobenius-invariant eight-arc in $\PG(2,25)$ admits an extension by a
fresh nonfixed point together with its Frobenius conjugate.
\end{theorem}

\begin{proof}
The fixed-point count is even and at most eight.  Proposition
\ref{prop:four-profiles} treats $f=0,4,6,8$, and Proposition~\ref{prop:f2}
treats $f=2$.
\end{proof}

The complete-arc classifications in
\cite{MarcuginiMilaniPambianco2007,CoolsaetSticker2009} show that the
smallest complete arcs in $\PG(2,25)$ have size $12$.  Thus every eight-arc
has an ordinary one-point extension.  That fact does not imply
Theorem~\ref{thm:q25}: ordinary extendability neither makes the legal-point
set Frobenius-stable nor ensures that a legal point and its conjugate are
jointly legal.

\begin{remark}[External census]\label{rem:census}
An external enumeration after fixed-point normalization reports $469{,}600$
invariant eight-arcs with profile $(f,e)=(2,3)$ and observed minimum legal
pair count $32$.  Independent C++ and Python implementations agree on the
census, minimum, minimizing orbit indices, and coordinate witness.  These
numbers are computational data, not inputs to Proposition~\ref{prop:f2} or
Theorem~\ref{thm:q25}.
\end{remark}

\section{Equivariant saturation}

Call an invariant arc \emph{pair-saturated} if it has no legal conjugate-pair
extension.  This is weaker than ordinary completeness.

\begin{corollary}[Orbit-saturation bound]\label{cor:saturation}
If a $\phi$-invariant $k$-arc in $\PG(2,s^2)$ is pair-saturated, then
\begin{equation}\label{eq:saturation}
  k\ge 1+\left\lceil\sqrt{2s(s-1)}\right\rceil.
\end{equation}
In particular, for every prime power $s\ge7$, every invariant eight-arc in
$\PG(2,s^2)$ admits a conjugate-pair extension.
\end{corollary}

\begin{proof}
If every fixed line is occupied, Lemma~\ref{lem:empty-lines} gives
$s^2+s+1=f(s+1)-\binom f2+e$.  Substituting $2e=k-f$ and completing the
square gives the occupation identity
\[
  k=s^2+1+(f-s-1)^2,
\]
which is stronger than \eqref{eq:saturation}.  Otherwise choose an empty
fixed line.  Pair saturation forces all $N=s(s-1)/2$ of its candidates to
be forbidden, so $N\le M$.  From $k=f+2e$ and \eqref{eq:M},
\[
  M=e((k-1)-e),
  \qquad
  4M\le(k-1)^2.
\]
Consequently $2s(s-1)\le(k-1)^2$, which gives
\eqref{eq:saturation}.

For $k=8$, one has $M\le12<N$ when $s\ge7$.  Full occupation is impossible
because $8<s^2+1$, so Theorem~\ref{thm:pair-criterion} applies.
\end{proof}

\begin{corollary}[All base orders from five]\label{cor:all-orders}
Let $s\ge5$ be a prime power.  Every Frobenius-invariant eight-arc in
$\PG(2,s^2)$ admits an extension by a fresh nonfixed point together with
its Frobenius conjugate.
\end{corollary}

\begin{proof}
The case $s=5$ is Theorem~\ref{thm:q25}.  There is no prime power strictly
between $5$ and $7$, and the final assertion of
Corollary~\ref{cor:saturation} treats every $s\ge7$.
\end{proof}

The scale in \eqref{eq:saturation} is the classical Lunelli--Sce scale for
ordinary complete arcs, with the same secant-covering mechanism in the
Frobenius quotient.  Ng and Wild's line-covering problem
\cite{NgWild2001} is another adjacent square-root-scale result.  The point
of Corollary~\ref{cor:saturation} is the weaker pair-saturation hypothesis,
not a new constant or a sharp family.

\section{Formal verification}\label{sec:formalization}

The results were formalized in Lean~\cite{Lean2015} on top of
mathlib~\cite{Mathlib2020}.  The supplementary source is divided as follows.

\begin{enumerate}[leftmargin=2.2em]
\item The projective-conjugation and quadratic-Frobenius modules construct
      the semilinear projective involution, prove the degree-two fixed-locus
      normalization, and count the $N=(s^2-s)/2$ candidates on a fixed line.
\item The line-counting and forbidden-candidate modules prove
      Lemma~\ref{lem:empty-lines}, the exact secant-orbit count, the
      injective charge, and the semantic arc-extension conclusion in
      Theorem~\ref{thm:pair-criterion}.
\item The global-count module defines the finite set of fresh nonfixed
      Frobenius pairs whose union with the old arc remains an arc.  It proves
      that this semantic set is the disjoint union of the carrierwise legal
      sets and that its cardinality is the abstract legal count.
\item The collision and invisibility modules prove
      Proposition~\ref{prop:linewise}, Corollary~\ref{cor:aggregate}, the
      fixed-center interpretation, and the cross-pair estimate
      \eqref{eq:cross-invisible}.
\item Separate profile modules prove the lower bounds $5$ and $4$ for
      $f=0$ and $f=4$.  The exceptional certificate proves
      Proposition~\ref{prop:f2}; the aggregate module proves
      Theorem~\ref{thm:q25} by the exhaustive parity split.
\item The orbit-saturation module proves the denominator-free inequality
      $2s(s-1)\le(k-1)^2$ used in Corollary~\ref{cor:saturation}.  The final
      ceiling notation and the geometric case split are paper proofs.
\end{enumerate}

The scoped builds and declaration-level axiom audits report exactly
\[
  [\texttt{propext},\ \texttt{Classical.choice},\ \texttt{Quot.sound}],
\]
the standard classical quotient profile inherited from mathlib.  No public
theorem depends on an admitted proof, a custom axiom, or
\texttt{native\_decide}.  The finite leaves in Proposition~\ref{prop:f2}
use kernel \texttt{decide}; the external census in Remark~\ref{rem:census}
is outside the theorem dependency graph.

For reproducibility, the paper-facing declarations are located as follows.

\begin{center}
\small
\begin{tabular}{>{\raggedright\arraybackslash}p{0.43\linewidth}
                >{\raggedright\arraybackslash}p{0.49\linewidth}}
\toprule
Source module & Paper-facing declaration or role \\
\midrule
\path{FiniteGeom/BaerCompletion/PairExtension.lean} &
  \path{quadraticBaer_pairExtension_lowerBound} \\
\path{RelativeConicArcs/QuadraticForbidden.lean} &
  \path{exists_quadratic_pair_extension} \\
\path{RelativeConicArcs/QuadraticGlobalCount.lean} &
  \path{card_globalLegalPairs_eq_legalCount} \\
\path{RelativeConicArcs/QuadraticInvisible.lean} &
  aggregate equality/excess and fixed-center criteria \\
\path{RelativeConicArcs/Q25ProfileZero.lean} &
  \path{five_le_sum_card_legal_profile_zero} \\
\path{RelativeConicArcs/Q25ProfileFour.lean} &
  \path{four_le_sum_card_legal_profile_four} \\
\path{RelativeConicArcs/Q25PairResult.lean} &
  \path{f2_pair_extension} \\
\path{RelativeConicArcs/Q25AllProfiles.lean} &
  \path{pair_extension} \\
\path{FiniteGeom/BaerCompletion/OrbitSaturation.lean} &
  \path{orbitSaturation_quadratic_bound_of_split} \\
\bottomrule
\end{tabular}
\end{center}

\subsection*{Artifact and reproduction}

The accompanying source archive pins Lean to
\path{leanprover/lean4:v4.32.0-rc1} and mathlib to revision
\path{571b8a8e54219b4d393f75f4b8653fac08197fcc}.  The audited formal-source
checkpoint is Git revision
\path{8f308c3d94652c3d3a66335701f1a18b53740675}.  From the \path{lean/}
directory the two paper-facing aggregate targets are rebuilt by
\begin{verbatim}
LEAN_NUM_THREADS=6 nix develop --command lake build \
  RelativeConicArcs.QuadraticGlobalCount \
  RelativeConicArcs.Q25AllProfiles
\end{verbatim}
The manuscript PDF is reproduced from its source directory with
\begin{verbatim}
nix shell nixpkgs#tectonic -c tectonic \
  frobenius_pair_extension.tex --keep-logs
\end{verbatim}

\section{Conclusion}

The fixed mate line of a quadratic Frobenius orbit turns symmetry-preserving
arc extension into a carrierwise counting problem.  The uniform count is a
first-order obstruction bound; its exact correction records which secant
orbits disappear at fixed centers and which visible obstructions collide.
This suffices, together with one finite exceptional-profile certificate, to
prove a uniform conjugate-pair extension theorem for invariant eight-arcs in
$\PG(2,25)$.

The natural strengthening is now family-specific rather than generic: prove
that every invariant eight-arc has at least two legal pairs, which would
give an alternate-orbit repair theorem for invariant ten-arcs after deletion
of one selected nonfixed orbit.  That multiplicity statement is not used or
claimed here.

\bibliographystyle{plain}
\bibliography{refs}

\end{document}
